\documentclass[10pt,a4paper]{article}
\usepackage[a4paper,left=14mm,right=14mm,top=18mm,bottom=15mm,headheight=19pt]{geometry}
\usepackage[T1]{fontenc}
\usepackage[utf8]{inputenc}
\usepackage[scaled=0.94]{helvet}
\renewcommand{\familydefault}{\sfdefault}
\usepackage{xcolor}
\usepackage{tabularx}
\usepackage{booktabs}
\usepackage{array}
\usepackage{fancyhdr}
\usepackage{lastpage}
\usepackage{tcolorbox}
\tcbuselibrary{skins}
\usepackage{hyperref}

% Clinical Color Palette
\definecolor{Navy}{HTML}{16324F}
\definecolor{Teal}{HTML}{0E7490}
\definecolor{SlateBg}{HTML}{F3F6FA}
\definecolor{WarnBg}{HTML}{FFFBEB}
\definecolor{WarnBorder}{HTML}{D97706}
\definecolor{FieldBg}{HTML}{F8FAFC}
\definecolor{BorderGray}{HTML}{B6C2CF}
\definecolor{Muted}{HTML}{475569}
\definecolor{RuleGray}{HTML}{CBD5E1}

\hypersetup{
  colorlinks=true,
  urlcolor=Teal,
  linkcolor=Navy,
  pdftitle={General CT Scan Imaging Forms — Fillable Clinical Workbook},
  pdfauthor={Medical Imaging Workbook}
}

% Header & Footer
\pagestyle{fancy}
\fancyhf{}
\fancyhead[L]{\fontsize{8}{10}\selectfont\bfseries\color{Navy} GENERAL CT SCAN IMAGING FORMS}
\fancyhead[R]{\fontsize{8}{10}\selectfont\color{Muted} Computed Tomography Clinical Workflow}
\fancyfoot[L]{\fontsize{7.5}{9}\selectfont\color{Muted} De-identified clinical \& educational use}
\fancyfoot[R]{\fontsize{7.5}{9}\selectfont\color{Muted} Page \thepage\ of \pageref{LastPage}}
\renewcommand{\headrulewidth}{0.5pt}
\renewcommand{\footrulewidth}{0.4pt}
\renewcommand{\headrule}{\hbox to\headwidth{\color{RuleGray}\leaders\hrule height \headrulewidth\hfill}}
\renewcommand{\footrule}{\hbox to\headwidth{\color{RuleGray}\leaders\hrule height \footrulewidth\hfill}}

\setlength{\parindent}{0pt}
\setlength{\parskip}{3pt}

% Interactive Form Field Helpers
\newcommand{\FormLabel}[1]{{\fontsize{7.8}{9.5}\selectfont\bfseries\color{Navy}#1}}

\newcommand{\FText}[3][15pt]{%
  \TextField[name=#2, width=#3, height=#1, bordercolor=BorderGray, backgroundcolor=FieldBg, charsize=8.5pt, borderwidth=0.6pt]{}%
}

\newcommand{\FMText}[3][30pt]{%
  \TextField[name=#2, width=#3, height=#1, multiline=true, bordercolor=BorderGray, backgroundcolor=FieldBg, charsize=8.5pt, borderwidth=0.6pt]{}%
}

\newcommand{\FCheck}[2]{%
  \CheckBox[name=#1, width=9pt, height=9pt, bordercolor=Navy, backgroundcolor=FieldBg, borderwidth=0.7pt]{}~{\fontsize{8}{10}\selectfont #2}%
}

\newcommand{\FieldRow}[2]{%
  \FormLabel{#1}\par\vspace{1pt}%
  \FText[15.5pt]{#2}{\linewidth}\par\vspace{3.5pt}%
}

\newcommand{\MFieldRow}[3][28pt]{%
  \FormLabel{#2}\par\vspace{1pt}%
  \FMText[#1]{#3}{\linewidth}\par\vspace{3.5pt}%
}

\newcommand{\PageBanner}[3]{%
  \noindent\begin{tcolorbox}[
    enhanced,
    colback=Navy,
    colframe=Navy,
    arc=2pt,
    boxrule=0pt,
    left=7pt,right=7pt,top=4.5pt,bottom=4.5pt,
    after skip=4pt
  ]
    {\fontsize{11.5}{14}\selectfont\bfseries\color{white}#1}
    \hfill
    \tcbox[on line, arc=2pt, colback=Teal, colframe=Teal, boxrule=0pt, left=5pt, right=5pt, top=1.5pt, bottom=1.5pt]{%
      \fontsize{7.5}{9}\selectfont\bfseries\color{white}#2%
    }
    \quad
    {\fontsize{8.5}{10}\selectfont\color{white!85!Navy} #3}
  \end{tcolorbox}%
}

\newtcolorbox{GuidanceBox}[1][]{
  enhanced,
  colback=SlateBg,
  colframe=RuleGray,
  arc=1.5pt,
  boxrule=0.5pt,
  left=6pt,right=6pt,top=4pt,bottom=4pt,
  after skip=4pt,
  #1
}

\newcommand{\RefBar}[1]{%
  \vfill
  \noindent\begin{tcolorbox}[
    enhanced,
    colback=SlateBg,
    colframe=RuleGray,
    arc=1.5pt,
    boxrule=0.4pt,
    left=6pt,right=6pt,top=3pt,bottom=3pt,
    after skip=0pt
  ]
    {\fontsize{7.5}{9.5}\selectfont\bfseries\color{Muted} References:\quad #1}
  \end{tcolorbox}%
}

\begin{document}
\begin{Form}

% ============================================================================
% PAGE 1: GENERAL CT REFERRAL & PRE-SCAN JUSTIFICATION
% ============================================================================
\PageBanner{General CT Scan (Computed Tomography)}{1. BEFORE THE SCAN}{Form 1 of 4}

\begin{GuidanceBox}
{\fontsize{8.1}{10.1}\selectfont
\textbf{\color{Navy}What to check first:} Confirm anatomical region, clinical question, prior imaging, and whether IV/oral contrast is needed.\par\vspace{1.5pt}
\textbf{\color{Navy}When to proceed or defer:} Use CT when indicated by clinical rules or urgent findings and when results will change care. If a validated low-risk rule or prior test answers the question, document deferral.\par\vspace{1.5pt}
\textbf{\color{Navy}Contrast \& radiation safety:} Review kidney function (eGFR/creatinine) when indicated, prior iodinated contrast reactions, pregnancy status, and pediatric/low-dose size-adjusted protocols.}
\end{GuidanceBox}

\noindent\begin{tabularx}{\linewidth}{@{}X@{\hspace{6pt}}p{2.4cm}@{\hspace{6pt}}p{2.6cm}@{\hspace{6pt}}X@{}}
\FormLabel{Patient ID} & \FormLabel{Age / DOB} & \FormLabel{Date} & \FormLabel{Ordering clinician} \\
\FText[15pt]{ct_p1_pid}{\linewidth} & \FText[15pt]{ct_p1_dob}{\linewidth} & \FText[15pt]{ct_p1_date}{\linewidth} & \FText[15pt]{ct_p1_clin}{\linewidth} \\
\end{tabularx}\par\vspace{3pt}

\noindent\begin{tabularx}{\linewidth}{@{}X@{\hspace{6pt}}p{4.6cm}@{}}
\FormLabel{Body region / CT protocol requested} & \FormLabel{Latest eGFR / Creatinine \& date} \\
\FText[15pt]{ct_p1_region}{\linewidth} & \FText[15pt]{ct_p1_egfr}{\linewidth} \\
\end{tabularx}\par\vspace{3pt}

\noindent\begin{tcolorbox}[enhanced,colback=SlateBg,colframe=RuleGray,arc=1.5pt,boxrule=0.5pt,left=6pt,right=6pt,top=3.5pt,bottom=3.5pt,after skip=3.5pt]
\FormLabel{Contrast:}\quad
\FCheck{ct_p1_cnc}{Noncontrast}\quad \FCheck{ct_p1_civ}{IV contrast}\quad \FCheck{ct_p1_cor}{Oral/enteric}\quad \FCheck{ct_p1_ccta}{CTA / Multiphasic}\qquad
\FormLabel{Technique:}\quad
\FCheck{ct_p1_tstd}{Standard}\quad \FCheck{ct_p1_tlow}{Low-dose}\quad \FCheck{ct_p1_tped}{Pediatric}\par\vspace{2.5pt}
\FormLabel{Prior contrast reaction:}\quad
\FCheck{ct_p1_rno}{None}\quad \FCheck{ct_p1_rmild}{Mild}\quad \FCheck{ct_p1_rsev}{Mod/Severe}\qquad
\FormLabel{Pregnancy:}\quad
\FCheck{ct_p1_pna}{N/A}\quad \FCheck{ct_p1_pneg}{Excluded}\quad \FCheck{ct_p1_ppos}{Pregnant (justified)}\quad \FCheck{ct_p1_punc}{Uncertain}
\end{tcolorbox}

\FieldRow{Clinical question to be answered}{ct_p1_cquest}

\MFieldRow[24pt]{History received (exact account, source \& relationship, recipient, date \& time)}{ct_p1_hist}

\MFieldRow[24pt]{Clinician assessment and prior tests / imaging (with dates and results)}{ct_p1_assess}

\noindent\begin{tabularx}{\linewidth}{@{}X@{\hspace{6pt}}X@{}}
\FormLabel{Guideline and exact scenario matched} & \FormLabel{Finding supporting ordering or deferral} \\
\FText[15pt]{ct_p1_guide}{\linewidth} & \FText[15pt]{ct_p1_find}{\linewidth} \\
\end{tabularx}\par\vspace{3pt}

\FieldRow{Patient-specific reason for an exception (if applicable)}{ct_p1_excep}

\noindent\begin{tcolorbox}[enhanced,colback=SlateBg,colframe=RuleGray,arc=1.5pt,boxrule=0.5pt,left=6pt,right=6pt,top=3.5pt,bottom=3.5pt,after skip=3.5pt]
\FormLabel{Decision:}\quad
\FCheck{ct_p1_ord}{Order}\quad \FCheck{ct_p1_def}{Defer}\qquad
\FormLabel{Urgency:}\quad
\FCheck{ct_p1_urout}{Routine}\quad \FCheck{ct_p1_uurg}{Urgent}\quad \FCheck{ct_p1_ustat}{Stat / Emergent}\par\vspace{2.5pt}
\FormLabel{What the result would change}\par\vspace{1pt}
\FText[14.5pt]{ct_p1_change}{\linewidth}
\end{tcolorbox}

\MFieldRow[22pt]{Safety decision and plain-language explanation for the patient}{ct_p1_safety}

\noindent\begin{tcolorbox}[enhanced,colback=SlateBg,colframe=RuleGray,arc=1.5pt,boxrule=0.5pt,left=6pt,right=6pt,top=3.5pt,bottom=3.5pt,after skip=2pt]
\FormLabel{Basis for request:}\quad
\FCheck{ct_p1_bclin}{Supported clinical use}\quad
\FCheck{ct_p1_bunc}{Uncertain clinical benefit}\quad
\FCheck{ct_p1_bres}{Research}\quad
\FCheck{ct_p1_bclar}{Awaiting clarification}\par\vspace{2.5pt}
\FormLabel{If research, study reference \& applicable approvals (or clarification note)}\par\vspace{1pt}
\FText[14.5pt]{ct_p1_resref}{\linewidth}
\end{tcolorbox}

\RefBar{\href{https://www.acr.org/Clinical-Resources/ACR-Appropriateness-Criteria}{ACR Appropriateness Criteria} \quad|\quad \href{https://www.acr.org/clinical-resources/clinical-tools-and-reference/contrast-manual}{ACR Contrast Manual} \quad|\quad \href{https://pubmed.ncbi.nlm.nih.gov/28937575/}{Pregnancy \& lactation}}

\newpage

% ============================================================================
% PAGE 2: GENERAL CT SCAN-SPECIFIC REPORT
% ============================================================================
\PageBanner{General CT Scan (Computed Tomography)}{2. SCAN-SPECIFIC REPORT}{Form 2 of 4}

\begin{GuidanceBox}
{\fontsize{8.2}{10.2}\selectfont\textbf{\color{Navy}What this report should record:} Record exact CT protocol, contrast agent/phases, coverage, dose metric if tracked, image quality/artefacts, pertinent positive and negative findings with measurements (mm/HU), incidental findings, and a concise impression answering the clinical question.}
\end{GuidanceBox}

\noindent\begin{tabularx}{\linewidth}{@{}X@{\hspace{6pt}}p{2.8cm}@{\hspace{6pt}}X@{}}
\FormLabel{Patient ID} & \FormLabel{Exam date} & \FormLabel{Reporting clinician} \\
\FText[15pt]{ct_p2_pid}{\linewidth} & \FText[15pt]{ct_p2_date}{\linewidth} & \FText[15pt]{ct_p2_rep}{\linewidth} \\
\end{tabularx}\par\vspace{3.5pt}

\noindent\begin{tabularx}{\linewidth}{@{}X@{\hspace{6pt}}X@{}}
\FormLabel{CT protocol, phases \& contrast (agent, volume, route)} & \FormLabel{Coverage, side / site \& prior comparisons} \\
\FText[15.5pt]{ct_p2_exam}{\linewidth} & \FText[15.5pt]{ct_p2_comp}{\linewidth} \\
\end{tabularx}\par\vspace{3.5pt}

\FieldRow{Clinical indication / question}{ct_p2_ind}

\noindent\begin{tcolorbox}[enhanced,colback=SlateBg,colframe=RuleGray,arc=1.5pt,boxrule=0.5pt,left=6pt,right=6pt,top=3.5pt,bottom=3.5pt,after skip=3.5pt]
\FormLabel{Technical adequacy:}\quad
\FCheck{ct_p2_tadeq}{Adequate}\quad
\FCheck{ct_p2_tsub}{Suboptimal}\quad
\FCheck{ct_p2_trep}{Repeat / additional phase needed}\qquad
\FormLabel{CTDIvol / DLP (optional):}\ \FText[13.5pt]{ct_p2_dose}{3.4cm}\par\vspace{2.5pt}
\FormLabel{Artefacts (motion, beam hardening, contrast timing) or coverage limitations}\par\vspace{1pt}
\FText[15pt]{ct_p2_qual}{\linewidth}
\end{tcolorbox}

\MFieldRow[82pt]{Findings and relevant measurements (organs, vessels, lymph nodes, bones, soft tissues; size in mm, attenuation in HU)}{ct_p2_find}

\MFieldRow[50pt]{Impression (direct answer to the clinical question and key incidental findings)}{ct_p2_imp}

\noindent\begin{tcolorbox}[enhanced,colback=SlateBg,colframe=RuleGray,arc=1.5pt,boxrule=0.5pt,left=6pt,right=6pt,top=3.5pt,bottom=3.5pt,after skip=3.5pt]
\FormLabel{Question status:}\qquad
\FCheck{ct_p2_qans}{Answered}\qquad
\FCheck{ct_p2_qpart}{Partly answered}\qquad
\FCheck{ct_p2_qinc}{Inconclusive}\par\vspace{2.5pt}
\FormLabel{Recommended next step:}\quad
\FCheck{ct_p2_nnone}{None}\quad
\FCheck{ct_p2_nct}{Follow-up CT}\quad
\FCheck{ct_p2_nmri}{MRI}\quad
\FCheck{ct_p2_nus}{Ultrasound}\quad
\FCheck{ct_p2_npet}{PET/CT}\quad
\FCheck{ct_p2_nbx}{Biopsy / Clinical corr.}
\end{tcolorbox}

\noindent\begin{tabularx}{\linewidth}{@{}X@{\hspace{6pt}}p{4.8cm}@{\hspace{6pt}}p{2.8cm}@{}}
\FormLabel{Follow-up action} & \FormLabel{Responsible clinician or team} & \FormLabel{Due date} \\
\FText[15.5pt]{ct_p2_fu}{\linewidth} & \FText[15.5pt]{ct_p2_resp}{\linewidth} & \FText[15.5pt]{ct_p2_due}{\linewidth} \\
\end{tabularx}\par\vspace{3.5pt}

\noindent\begin{tabularx}{\linewidth}{@{}X@{\hspace{6pt}}p{4.0cm}@{\hspace{6pt}}p{4.2cm}@{}}
\FormLabel{Urgent communication: recipient} & \FormLabel{Date \& time} & \FormLabel{Method} \\
\FText[15.5pt]{ct_p2_uwho}{\linewidth} & \FText[15.5pt]{ct_p2_utime}{\linewidth} & \FText[15.5pt]{ct_p2_umeth}{\linewidth} \\
\end{tabularx}

\RefBar{\href{https://www.acr.org/Clinical-Resources/ACR-Appropriateness-Criteria}{ACR Appropriateness Criteria} \quad|\quad \href{https://www.rsna.org/practice-tools/data-tools-and-standards/radreport-reporting-templates}{RSNA RadReport Templates}}

\newpage

% ============================================================================
% PAGE 3: WHAT THE IMAGE SHOWS AND MEANS
% ============================================================================
\PageBanner{General CT Scan (Computed Tomography)}{3. WHAT THE IMAGE SHOWS \& MEANS}{Form 3 of 4}

\renewcommand{\arraystretch}{1.2}
\noindent\begin{tabularx}{\linewidth}{@{}>{\bfseries\color{Navy}}p{3.0cm} >{\color{Navy}}p{5.4cm} X@{}}
\toprule
\textbf{Context} & \textbf{Key CT Finding} & \textbf{Typical Meaning / Interpretation} \\
\midrule
Head / Trauma & Hyperdense collection (50--80 HU); midline shift & Acute blood (extra-axial/parenchymal); mass effect \\
Chest / CTA & Intraluminal filling defect; focal consolidation & Pulmonary embolism/thrombus; infection/inflammation \\
Abdomen / Pelvis & Fat stranding, wall thickening, free air & Acute inflammation, obstruction, or perforation \\
Lesion (HU) & Attenuation ($<0$ fat; $\sim$0--20 fluid; $>$100 calcium) & Helps narrow tissue composition and enhancement \\
Oncology / Staging & Enlarged nodes, organ lesion, interval change & Supports staging, progression, or treatment response \\
\bottomrule
\end{tabularx}\par\vspace{4pt}

\begin{GuidanceBox}
{\fontsize{8.1}{10.1}\selectfont
\textbf{\color{Navy}What CT leaves unanswered:} A normal noncontrast CT does not exclude concussion, early ischaemic stroke, subtle marrow/ligament injury, or isoattenuating solid-organ lesions. Motion, metallic streak, or suboptimal contrast bolus timing can obscure findings.\par\vspace{2pt}
\textbf{\color{Navy}Worked interpretation example:} Finding: Hyperdense collection outside the brain. Interpretation: may indicate acute extra-axial blood. If absent, check the limitations above before closing the clinical question.\par\vspace{2pt}
\textbf{\color{Navy}Illustration guide (4:3 plate):} Label an illustrative axial/coronal/sagittal slice with window setting (e.g., brain, lung, soft-tissue, bone) and mark key measurements or midline shift.}
\end{GuidanceBox}

\noindent\begin{tabularx}{\linewidth}{@{}X@{\hspace{6pt}}p{4.0cm}@{\hspace{6pt}}p{2.6cm}@{\hspace{6pt}}p{2.0cm}@{}}
\FormLabel{Primary finding} & \FormLabel{Location / Side} & \FormLabel{Measurement} & \FormLabel{Unit (mm / HU)} \\
\FText[15.5pt]{ct_p3_find}{\linewidth} & \FText[15.5pt]{ct_p3_loc}{\linewidth} & \FText[15.5pt]{ct_p3_meas}{\linewidth} & \FText[15.5pt]{ct_p3_unit}{\linewidth} \\
\end{tabularx}\par\vspace{3.5pt}

\noindent\begin{tabularx}{\linewidth}{@{}X@{\hspace{6pt}}p{4.0cm}@{\hspace{6pt}}p{2.6cm}@{\hspace{6pt}}p{2.0cm}@{}}
\FormLabel{Secondary finding (if any)} & \FormLabel{Location / Side} & \FormLabel{Measurement} & \FormLabel{Unit (mm / HU)} \\
\FText[15.5pt]{ct_p3_find2}{\linewidth} & \FText[15.5pt]{ct_p3_loc2}{\linewidth} & \FText[15.5pt]{ct_p3_meas2}{\linewidth} & \FText[15.5pt]{ct_p3_unit2}{\linewidth} \\
\end{tabularx}\par\vspace{3.5pt}

\noindent\begin{tcolorbox}[enhanced,colback=SlateBg,colframe=RuleGray,arc=1.5pt,boxrule=0.5pt,left=6pt,right=6pt,top=3.5pt,bottom=3.5pt,after skip=3.5pt]
\FormLabel{Hypothesis impact:}\quad
\FCheck{ct_p3_hsup}{Supports hypothesis}\qquad
\FCheck{ct_p3_hless}{Makes it less likely}\qquad
\FCheck{ct_p3_hunres}{Leaves it unresolved}
\end{tcolorbox}

\MFieldRow[36pt]{Plain-language explanation for patient / care team}{ct_p3_plain}

\noindent\begin{tcolorbox}[enhanced,colback=SlateBg,colframe=RuleGray,arc=1.5pt,boxrule=0.5pt,left=6pt,right=6pt,top=4pt,bottom=4pt,after skip=2pt]
\FormLabel{Audit Challenge}\par\vspace{2pt}
\noindent\begin{tabularx}{\linewidth}{@{}p{5.5cm}@{\hspace{6pt}}X@{}}
\FormLabel{Reviewer \& date} & \FormLabel{Finding disputed} \\
\FText[15pt]{ct_p3_arev}{\linewidth} & \FText[15pt]{ct_p3_adisp}{\linewidth} \\
\end{tabularx}\par\vspace{2.5pt}
\MFieldRow[24pt]{Supporting evidence}{ct_p3_aevid}
\MFieldRow[24pt]{Response or resolution}{ct_p3_ares}
\end{tcolorbox}

\RefBar{\href{https://www.radiologyinfo.org/en/info/bodyct}{RadiologyInfo: Body CT} \quad|\quad \href{https://www.radiologyinfo.org/en/info/headct}{RadiologyInfo: Head CT}}

\newpage

% ============================================================================
% PAGE 4: FIND SIMILAR CASES AND CHECK THE EVIDENCE
% ============================================================================
\PageBanner{General CT Scan (Computed Tomography)}{4. SIMILAR CASES \& EVIDENCE}{Form 4 of 4}

\begin{GuidanceBox}[colback=WarnBg,colframe=WarnBorder!50!white,borderline west={3pt}{0pt}{WarnBorder}]
{\fontsize{8.1}{10.1}\selectfont\textbf{\color{Navy}Optional review aid:} Use de-identified clinical terms only. Keep patient names, identifiers, reports, and patient images out of public searches. Record examples for review without delaying care.}
\end{GuidanceBox}

\begin{GuidanceBox}
{\fontsize{8.1}{10.2}\selectfont
\textbf{\color{Navy}1. Check the decision using only pre-scan information (example templates):}\par
\texttt{<clinical presentation> <anatomic region> CT appropriateness guideline}\par\vspace{3pt}
\textbf{\color{Navy}2. Find a teaching case using the imaging findings (example templates):}\par
\texttt{<contrast/noncontrast> <region> CT <key finding> <measurement/pattern> teaching case}}
\end{GuidanceBox}

\FieldRow{Actual search terms used}{ct_p4_terms}

\noindent\begin{tabularx}{\linewidth}{@{}X@{\hspace{6pt}}p{3.4cm}@{}}
\FormLabel{Source link / citation} & \FormLabel{Access date} \\
\FText[15.5pt]{ct_p4_src}{\linewidth} & \FText[15.5pt]{ct_p4_date}{\linewidth} \\
\end{tabularx}\par\vspace{5pt}

\MFieldRow[62pt]{What matches and what differs}{ct_p4_match}

\MFieldRow[62pt]{How the diagnosis was established (or whether confirmation was unclear)}{ct_p4_diag}

\MFieldRow[62pt]{What the source adds and what remains uncertain}{ct_p4_adds}

\RefBar{\href{https://www.acr.org/Clinical-Resources/ACR-Appropriateness-Criteria}{ACR Appropriateness Criteria} \quad|\quad \href{https://radiopaedia.org/}{Radiopaedia (CT Teaching Cases)}}

\end{Form}
\end{document}
